\labsection{8}{Genetic Algorithm optimization}{sec:lab08-ga}

\begin{labproject}{Five GA formulations in MATLAB}
\textbf{Coverage.} Chapter~\ref{ch:ga} --- five GA formulations.\\
\textbf{Source files.} \texttt{Lab08/CODE\_1.m} through \texttt{CODE\_5.m}, one file per problem.\\
\textbf{Goal.} Keep the GA loop recognizable while changing chromosome representation and fitness definition to match each problem.

\begin{labmode}
Run the five baseline loops first; stochastic initialization and operators may produce different results. Each revised script reports the best final chromosome in problem-specific terms, not only a convergence curve.
\end{labmode}

\begin{pitfall}
Without elitism, later variation can remove the best chromosome. Distinguish the best fitness \emph{seen} from the best individual that \emph{survives} in the final population.
\end{pitfall}

\medskip\noindent\textbf{Part A --- Continuous one- and two-variable objectives.}\par
Run Code 1 and Code 2. Plot best fitness versus generation and record the best chromosome found at the end of the run.

\medskip\noindent\textbf{Part B --- Linear regression.}\par
Run Code 3 and report the best discovered slope/intercept together with the final MSE. Compare the values qualitatively with the data-generating parameters \([2.5,5]\).

\medskip\noindent\textbf{Part C --- Binary knapsack.}\par
Run Code 4. Decode the best chromosome into selected items, total weight, and total value. Confirm that the capacity constraint is satisfied.

Because only $2^4=16$ selections exist, compare the GA result with the brute-force optimum.

\begin{keyidea}
An over-capacity chromosome receives fitness 0: a \term{penalty}, not a repair. Repairing or rejecting it would produce different search behavior.
\end{keyidea}

\medskip\noindent\textbf{Part D --- Polynomial roots.}\par
Run Code 5 and report whether the final population converges nearer to \(-2\) or 3 in that run. Repeat with a different random seed and compare.

\begin{tryit}
Add elitism to one code by copying the current best chromosome unchanged into the offspring. Compare the original and elitist traces and final solutions.
\end{tryit}

\begin{troubleshoot}
\textbf{Minimization.} Roulette selection requires larger scores for better candidates. Codes 2 and 5 transform the objective during selection:

\begin{itemize}
  \item Code 2 keeps the raw Rosenbrock value, then selects on \texttt{invFitness = 1 ./ (fitness + 1e-6)}, so a smaller objective gives a larger share of the wheel.
  \item Code 5 negates the objective, then shifts with \texttt{prob = fitness - min(fitness) + eps} so the worst individual sits at zero probability rather than a negative one.
\end{itemize}

Both transformations assign larger selection scores to better candidates.

\textbf{Run-to-run variation.} GAs are stochastic; report results across several runs.

\textbf{Different roots in Code 5.} Both $-2$ and $3$ are global minima with $f=0$.

\textbf{Suboptimal knapsack result.} Increase population or generations, or add elitism.
\end{troubleshoot}
\end{labproject}

\begin{verification}
\begin{itemize}
  \item Each script reports its decoded solution and convergence plot.
  \item The knapsack chromosome is decoded into items, weight, and value, and the capacity constraint is checked.
  \item Code 3's parameters are compared with $[2.5,5]$.
  \item Code 5 is run more than once, with the root reported for each run.
  \item Differences between peak and final fitness are explained.
\end{itemize}
\end{verification}

\begin{labdeliverable}
Submit one MATLAB live script containing all five GA problems, five convergence plots, the final decoded solution for each problem, and a comparison table listing chromosome type, fitness direction, crossover type, and mutation type.
\end{labdeliverable}
